\documentclass{article}
\usepackage[includefoot]{geometry}
\usepackage{/home/user1/code/tex/sty}
\usepackage{amsmath}
\usepackage{floatrow}
\usepackage{hyperref}
\geometry{top=0.75in, bottom=1in, left=0.75in, right=0.75in}
\usepackage[backend=bibtex]{biblatex}
\addbibresource{bib.bib}
\floatsetup[figure]{objectset=centering}
\floatsetup[subfigure]{objectset=centering}
\hypersetup{colorlinks=true,citecolor=black,linkcolor=black}
\setlength{\parindent}{0pt}
\title{Paper}
\author{Agustin Romero\thanks{ar1@stanford.edu}}
\newcommand{\figpic}[3]{
\begin{figure}[H]
\begin{center}
\pic{#1}{#2}
\caption{#3}
\end{center}
\end{figure}}
\newcommand{\figtab}[3]{
\begin{table}[H]
\begin{tabular}{#1}
#2
\end{tabular}
\caption{#3}
\end{table}}
\begin{document}
\maketitle
This paper studies the renormalization group flow of Fermi Liquid Theory. The Hamiltonian for the fermi liquid will be described as
\g{H=V\int d^3\vec{k} (\epsilon_k - \mu) \psi^\dagger_k \psi_k + \br{1}{2}V\int f(k_1,k_2,q) \psi^\dagger_{k_1+q} \psi_{k_2-q} \psi_{k_2} \psi_{k_1}  d^3k_1 d^3 k_2 d^3 q}
The Fermi surface is identifiable at 
\g{E(\vec{k})=\mu}
and the existence of only one band is assumed. The procedure for analyzing the Fermi liquid Hamiltonian will be as follows. Define a partition function $Z$ in the limit $T=0$. Then we will define the renormalization group transformations for the coupling and fields such that the free interaction term is invariant, obtaining a free field fixed point in the RG flow. Then we will enumerate the terms around the fixed point, beginning with tree level interactions and then loop-level interactions. If an interaction is relevant, then the excitations will have a gap. Conversely, if no relevant operators appear, then the excitations will be gapless from the Fermi surface.\\~\\
We consider a weak scattering process $k_1 + k_2 \rightarrow k_3 + k_4$ between energy states that are close to the Fermi surface. There are three primary ways of considering the momenta. 
\begin{enumerate}
\item Each of $k_1$, $k_2$, $k_3$ and $k_4$ must lie exactly on the Fermi sphere. This is not typically the assumption made, however it does illustrate that the direction of $\vec{q}$ is highly restricted: $\vec{q}$ must point tangential the Fermi sphere.
\item $\vec{q}=0$. This is forward scattering and is denoted by $f$.
\item $\vec{k_1}=-\vec{k_2}$. This is the ,,Cooper channel'' or ,,BCS'' channel and is denoted by $g$.
\end{enumerate}
The interaction potential $f(k_1,k_2,q)$ can be split up into multiple cases, in particular, one for the forward scattering case and one for the Cooper channel case.
\g{H_p = \int d^3 \vec{k} d^3\vec{k'} f_{kk'} n_k n_{k'} + \int d^3{k}d^3{k'} g_{kk'} \psi^\dagger_{-k'} \psi_{k'} \psi_k \psi_{-k}}
We would like to study the $\beta$ function flow of $f$ and $g$. We will find that the different cases of momenta correspond to different $\beta$ flows.
\g{||\vec{k_i}-k_F|\leq \Lambda}
By introducing the cutoff $\Lambda$, the loop momentum must be within a shell $d\Lambda$. Next the action is
\g{S_0=\sum_i \int_{-\Lambda}^{\Lambda} \int_{-\infty}^{\infty} \br{d\omega}{2\pi} \overline{\psi_i}(\omega k) (-i\omega -k)\psi_i(\omega k)}
The partition function is
\g{Z_0=\int \prod_i \prod_{|k|\leq \Lambda} d\overline{\psi}_i(\omega k) d\psi_i(\omega k) e^{S_0}}
Integrate $\psi(k\omega)$ over the range $\br{\Lambda}{s}\leq |\vec{k}| \leq \Lambda$ to obtain a prefactor. This prefactor is thrown out, and the remaining momenta in the interaction are restricted to the range
\g{-\br{\Lambda}{s}\leq |\vec{k}| \leq \br{\Lambda}{s}}
Then, rescale the variables.
\g{k'&=sk\\
\omega'&= s\omega \\
\psi'(k' \omega') &=s^{-3/2} \psi_i(k\omega)}
This rescaling of momentum and frequency of the interaction results in the modified action,
\g{S'=s^3\sum_i \int_{-\Lambda}^{\Lambda} \br{dk'}{2\pi}\int_{-\infty}{^\infty} \psi_i (\br{\omega'}{s},\br{k'}{s})(i\omega' - k')\psi_i(\br{\omega'}{s},\br{k'}{s}}
We want to rescale the momenta and fields such that the free action $S$ is initially a fixed point in the system, and then consider perturbations of interactions. We have already rescaled the variables, now in order to return to the original action $S$, we must now rescale the fields.
\g{\psi_i(\br{\omega'}{s},\br{k'}{s}}
Resulting in the action $S'=S$,
\g{S'=s^3\sum_i \int_{-\Lambda}^{\Lambda} \br{dk'}{2\pi}\int_{-\infty}^{\infty} \psi_i (\omega',k')(i\omega' - k')\psi_i(\omega',k')=S}
\subsection{Quadratic Perturbation}
\g{\delta S_2 = \sum_i \int_{-\Lambda}^{\Lambda} \br{dk}{2\pi} \int_{-\infty}{^\infty} \br{d\omega}{2\pi} \mu(k\omega) \overline{\psi_i}(k\omega)\psi_i(k\omega}
The rescaled Fermi surface is
\g{\mu'(\omega',k',i)=s\mu(\omega,k,i)}
which can be expanded in the Taylor series
\g{\mu(k,\omega)=\mu_{00}+\mu_{10}k+\mu_{01}i\omega+...+\mu_{nm}(i\omega)^m+...}
The $\mu_{00}$ term is a relevant perturbation,
\g{\mu_{00}\rightarrow s\mu_{00}}
This results in a rescaling of the chemical potential and the Fermi sea by the factor $s$.
\subsection{Quartic Perturbation}
Let the function $u$ determine the interaction at the quartic level. We are interested in evaluating the quartic perturbation to the action and to evaluate its RG flow. The quartic coupling is calculated relative to the origin $k=0$ and so
\g{K_i=\epsilon_i(K_F+k)}
represents the momentum state of a particular Fermi field $\psi(K_i)$ where $\epsilon_i=\pm 1$ for R,L respectively. Let $\psi(i)=\psi(K_i,\omega_i)$ for concision. The quartic coupling must be integrated over all four momenta $K_i$ and $\omega_i$, and so the integration is already quite complicated, in particular,
\g{\int_{K\omega}&=\prod_{i=1}^4 \int_{-\pi}^\pi \br{dK_i}{2\pi} \int_{-\infty}^\infty \br{d\omega_i}{2\pi} 2\pi \delta(K_1+K_2-K_3-K_4)2\pi\delta(\omega_1+\omega_2-\omega_3-\omega_4)
\intertext{Expand the $K_i$.}
&= \int_{-\Lambda}^\Lambda \br{dk_1...dk_4}{(2\pi)^4} \int_{-\infty}^\infty \br{d\omega_1...d\omega_4}{(2\pi)^4} 2\pi\delta(\omega_1+\omega_2-\omega_3-\omega_4) \\ &\times 2\pi \delta(\epsilon_{i_1}(K_f+k_1)+\epsilon_{i_2}(K_f+k_2)-\epsilon_{i_3}(K_f+k_3)-\epsilon_{i_4}(K_f+k_4))}
This integration $\int_{K\omega}$ can now be plugged into the quartic term,
\g{\delta S_4 &= \br{1}{2!2!} \int_{K\omega} \overline{\psi}(4)\overline{\psi}(3)\psi(2)\psi(1)u(4,3,2,1)\\
&= \br{1}{2!2!} \sum_{i_1,i_2,i_3,i_4} \int_{K\omega}^\Lambda \overline{\psi}_{i_4}(4)\overline{\psi}_{i_3}(3)\psi_{i_2}(2)\psi_{i_1}(1)u_{i_4,i_3,i_2,i_1}(4,3,2,1)}
This is the 4-point contribution to the action at tree level, where the real meaning of the interaction is represented by the function $u$. Depending on the scaling of $u$ we can determine the RG flow of this 4-point interaction. We can calculate its scaling by dimensional analysis, because we can compute the scaling of the other terms. In particular, we must compute the scaling of the $\delta$ functions and $\psi$ fields. First rescale the variables to primed variables, and then rescale the fields to the primed fields.
\g{\delta(\omega_1'+\omega_2'-\omega_3'-\omega_4')=s\delta(\br{\omega_1'}{s}+\br{\omega_2'}{s}-\br{\omega_3'}{s}-\br{\omega_4'}{s})}
The $\delta$ over $\omega$ gives a factor of $s$. The $\delta$ over $K$ is 
\g{\delta(K_1+K_2+K_3+K_4)=s\delta(K_f\sum_i \epsilon_i + \sum_i \epsilon_i k_i)}
This $\delta$ function enforces momentum conservation. We know very simply that $K_F$ should not enter into the momentum conservation at all, only the $K_i$. So, there are only certain allowed values of $i_4i_3i_2i_1$ due to momentum conservation, such that the $K_F$ term cancels. These are
\g{LRLR,RLRL,RLLR,LRRL}
resulting in either an overall phase factor, or 0,
\g{\sum_i\epsilon_i=0,\pm 4 \qquad \sum_i K_f\epsilon_i = 0,\pm 4\pi K_F}
The overall phase factor is permissible, because the $\delta$ function is periodic in $4\pi$. So ultimately we are safe to conclude that the $\delta$ over momenutum is independent of $K_F$, and the scaling is simply
\g{\delta(k)\rightarrow \delta(\br{k'}{s})}
and therefore the variable definition is simply
\g{\delta(\br{k'}{s})=s\delta'(k')}
The $\delta$ over $K$ gives a factor of $s$ as well. In total, the variable redefinitions give an overall factor, schematically,
\g{\prod \br{dk'}{s} \prod \br{d\omega'}{s} s \delta(k') s\delta(\omega')}
of $s^{-6}$. Now, we must rescale the fields according to
\g{\psi'=s^{3/2}{\psi}}
as done previously in the quadratic case. This provides exactly a factor of $s^6$. This reveals that the scaling of $u$ is of order $s^0$.
\g{u_{i_4i_3i_2i_1}(\br{\omega'}{s},\br{k_i'}{s})=u'_{i_4i_3i_2i_1}(\omega'_i,k'_i)}
\end{document}
